\chapter{Electrocardiogram (ECG)}
\index{Electrocardiogram (ECG)}

\section*{Definition and Overview}
An electrocardiogram, commonly called an ECG or EKG, is a quick and painless test that records the electrical activity of the heart. The heart's contractions are triggered by small electrical impulses, and electrodes placed on the skin detect these signals and display them as a tracing. The ECG shows the heart rate and rhythm, and it reveals evidence of heart attacks, enlarged heart chambers, and other cardiac conditions. It is one of the most commonly performed tests in medicine, used in GP clinics, emergency departments, and hospitals.

\section*{Epidemiology}
The ECG is used in the assessment of millions of people each year. It plays a central role in the rapid diagnosis of heart attacks, which remain a leading cause of death, and in the detection of abnormal heart rhythms, which affect a significant proportion of older adults. It is also used in routine health checks, pre-operative assessment, and the monitoring of people with known heart disease.

\section*{Aetiology and Pathophysiology}
The ECG traces the heart's electrical cycle.
\begin{itemize}
    \item \textbf{The heart's pacemaker:} the sinoatrial node generates each heartbeat's electrical impulse
    \item \textbf{Conduction:} the impulse spreads through the atria and then the ventricles through specialised pathways
    \item \textbf{The waves:} the P wave reflects atrial contraction, the QRS complex ventricular contraction, and the T wave ventricular recovery
    \item \textbf{What the ECG detects:} heart rate and rhythm, conduction delays, heart attack damage, chamber enlargement, and the effects of medicines and electrolyte imbalances
    \item \textbf{Interpretation:} the tracing is read for patterns that indicate normal or abnormal electrical activity
    \item \textbf{Limitations:} a normal ECG does not rule out all heart problems, and abnormalities may be intermittent
\end{itemize}

\section*{Clinical Features}
An ECG may be requested in many situations.
\begin{itemize}
    \item Chest pain, pressure, or tightness
    \item Palpitations, fluttering, or an irregular heartbeat
    \item Breathlessness, dizziness, or fainting
    \item Routine assessment for people with heart disease risk factors
    \item Pre-operative checks
    \item Monitoring after a heart attack or heart surgery
    \item Checking the effects of heart medicines
    \item Exercise stress testing
\end{itemize}

\section*{Differential Diagnosis}
The ECG result is interpreted in the context of the person's symptoms.
\begin{itemize}
    \item A resting ECG can be normal despite significant coronary disease
    \item Some ECG variants are harmless, while others are significant
    \item Comparison with previous tracings helps interpretation
    \item Blood tests and imaging complement the ECG
\end{itemize}

\section*{When to Seek Medical Care}
An ECG is performed when a clinician considers it appropriate. Symptoms such as chest pain, palpitations, breathlessness, or fainting should always be assessed by a doctor.

\textbf{Contact your local emergency services for:}
\begin{itemize}
    \item Chest pain lasting more than a few minutes, especially with sweating, breathlessness, or nausea
    \item Fainting or collapse
    \item Palpitations with dizziness or chest pain
    \item Stroke symptoms, including facial drooping, weakness, or difficulty speaking
\end{itemize}

\section*{Diagnosis and Investigations}
Undergoing an ECG is simple and comfortable.
\begin{itemize}
    \item \textbf{Preparation:} the person removes upper clothing and lies down; electrodes are placed on the chest, arms, and legs
    \item \textbf{Recording:} the test takes only a few minutes and is painless
    \item \textbf{Interpretation:} a doctor reads the tracing, often with computer assistance
    \item \textbf{Extended monitoring:} Holter monitors record over 24 hours for intermittent symptoms; exercise ECGs record during activity
    \item \textbf{Cardiac blood tests:} troponin measurements accompany the ECG when a heart attack is suspected
\end{itemize}

\section*{Management}
The ECG guides management of the underlying condition.
\begin{itemize}
    \item \textbf{Normal result:} reassurance, with further assessment if symptoms continue
    \item \textbf{Heart attack:} urgent treatment to restore blood flow
    \item \textbf{Arrhythmias:} medicines, cardioversion, or other procedures
    \item \textbf{Other findings:} management of the cause with cardiology referral as needed
    \item \textbf{Monitoring:} repeat ECG when symptoms change or during treatment
\end{itemize}

\section*{Complications}
\begin{itemize}
    \item The test itself is safe with no significant complications
    \item Skin irritation from electrodes in a small number of people
    \item The risk of misinterpretation, which is minimised by specialist reading
    \item False reassurance from a normal resting ECG
\end{itemize}

\section*{Prognosis and Outcomes}
The ECG is fast, inexpensive, and life-saving, enabling rapid diagnosis of heart attacks and dangerous rhythms. Used with other tests and clinical assessment, it supports effective treatment of a wide range of heart conditions and gives valuable reassurance when results are normal.

\section*{Prevention and Self-Care}
\begin{itemize}
    \item Seek prompt attention for chest pain, palpitations, or fainting
    \item Attend health checks if you have heart disease risk factors
    \item Know where your previous ECGs are stored
    \item Manage blood pressure, cholesterol, and other risk factors with your doctor
    \item Do not rely on consumer heart monitors for diagnosis
\end{itemize}

\section*{Sources and Further Reading}
The following reputable sources provide further information for healthcare students and patients seeking reliable, current advice about the electrocardiogram.
\begin{itemize}
    \item Heart Foundation Australia. \textit{Heart tests: ECG.} https://www.heartfoundation.org.au/
    \item Healthdirect Australia. \textit{Electrocardiogram (ECG).} https://www.healthdirect.gov.au/electrocardiogram
    \item NHS. \textit{Electrocardiogram (ECG).} https://www.nhs.uk/conditions/electrocardiogram/
    \item Mayo Clinic. \textit{Electrocardiogram (ECG or EKG).} https://www.mayoclinic.org/tests-procedures/ekg/
    \item MedlinePlus. \textit{Electrocardiogram.} https://medlineplus.gov/ency/article/003868.htm
    \item Better Health Channel. \textit{ECG tests.} https://www.betterhealth.vic.gov.au/
\end{itemize}
